\documentclass[main]{subfiles}
\usepackage[sols]{workbook}
\numbering{ws2-2}

\begin{document}

\ifSubfilesClassLoaded{%
  \chaptermark{\chaptertwo}
}{}

\section{Taylor Approximations II} \label{ws2-2}

\begin{framed}
\centerline{\textbf{Lesson Goals}}

You should be able to:
\begin{itemize}
\item Find the degree $n$ Taylor polynomial of a given function $f(x)$ centered at $x=c$, using the Taylor coefficient formula
\item Use a degree $n$ Taylor polynomial to approximate a specific function value, by identifying an appropriate function $f(x)$ and center $c$
\item Obtain information about a function's behavior near the center from the Taylor coefficients.
\item Express basic sums of terms in both $\sum$ notation and $\cdots$ notation
\end{itemize}
\end{framed}

\begin{enumerate}


  \item 
 The goal of Taylor approximation is to approximate a function $f(x)$ with an $n$th degree polynomial $P_n(x)$. To do so, we first pick a center value $x=c$ and decide on a degree $n$. Then, the polynomial will look like:
  $$\boxed{P_n(x)=a_0+a_1(x-c)+a_2(x-c)^2+ \cdots + a_{n-1}(x-c)^{n-1}+a_n(x-c)^n}$$
  where $a_0, a_1, a_2, ..., a_n$ are numerical coefficients to be found.
  
  To make the polynomial a good approximation of $f(x)$, we want its derivatives at $x=c$ to match the derivatives of $f$ at $x=c$. That is, we want all of the following equations to be true:
  \begin{align*}
  P_n(c) &= f(c)\\
  P_n'(c) &= f'(c)\\
  p_n''(c) &= f''(c)\\
  &~~\vdots \\
  p_n^{(n)}(c)&=f^{(n)}(c)
  \end{align*}
  Using these equations, we can solve for the coefficients $a_0, a_1, a_2, ..., a_n$ in general. 
  
For each part below, let $P_n(x)=a_0+a_1(x-c)+a_2(x-c)^2+ a_3(x-c)^3+a_4(x-c)^4+\cdots +a_n(x-c)^n$.
  
 \begin{enumerate}
 \item 
 \begin{problem}[\vfill]
 Find $P_n'(x)$, $P_n''(x)$, $P_n^{(3)}(x)$, and $P_n^{(4)}(x)$, the first four derivatives of $P_n(x)$. Write your answers in $\cdots$ notation. Don't simplify the coefficients! (Remember that $a_k$, $c$, and $n$ are all constants.)
 \end{problem}
 
 
 \begin{solution}
 \begin{align*}
 P_n'(x) &=a_1 + 2a_2(x-c)+3a_3(x-c)^2+4a_4(x-c)^3+\cdots + n\cdot a_n(x-c)^{n-1}\\
 P_n''(x) &= 2a_2 + 3\cdot 2a_3(x-c)+4\cdot 3(x-c)^2+\cdots + n\cdot (n-1) \cdot a_n(x-c)^{n-2}\\
P_n^{(3)}(x) &= 3\cdot 2a_3+4\cdot 3\cdot 2a_4(x-c)+\cdots + n\cdot (n-1) \cdot (n-2)\cdot a_n(x-c)^{n-3}\\
P_n^{(4)}(x) &= 4\cdot 3\cdot 2a_4+\cdots + n\cdot (n-1) \cdot (n-2)\cdot (n-3)\cdot a_n(x-c)^{n-4}
 \end{align*}
 \end{solution}
 
 

 \pspace{\newpage}

 
 \item 
 \begin{problem}[1in]
 Use the equation for $P_n(x)$, along with your results from part (a), to find $P_n(c)$, $P_n'(c)$, $P_n''(c)$, $P_n^{(3)}(c)$, and $P_n^{(4)}(c)$. 
 \end{problem}
 
 
 \begin{solution}
 Plugging $c$ in for $x$ in our previous equations from part (a), we get:
 \begin{align*}
 P_n(c)&=a_0\\
 P_n'(c)&=a_1\\
 P_n''(c)&=2a_2\\
 P_n^{(3)}(c)&=3\cdot2a_3\\
 P_n^{(4)}(c)&=4\cdot 3\cdot 2a_4
 \end{align*}
 \end{solution}
 
 
  Recall that $P_n(x)=a_0+a_1(x-c)+a_2(x-c)^2+ a_3(x-c)^3+a_4(x-c)^4+\cdots +a_n(x-c)^n$.
 \item 
 \begin{problem}[1.5in]
 Using your results from part (b), along with the equations
   \begin{align*}
  P_n(c) &= f(c)\\
  P_n'(c) &= f'(c)\\
  P_n''(c) &= f''(c)\\
  P_n^{(3)}(c)&=f^{(3)}(c)\\
  P_n^{(4)}(c)&=f^{(4)}(c)
  \end{align*}
  express $a_0$, $a_1$, $a_2$, $a_3$, and $a_4$ in terms of $f$. 
 \end{problem}
 
 \begin{solution}
 By plugging in our results from part (b) into the equations above, we get:
 
  \begin{align*}
  P_n(c) = f(c) &\Longrightarrow \boxed{a_0=f(c)}\\
  P_n'(c) = f'(c)&\Longrightarrow \boxed{a_1=f'(c)}\\
  P_n''(c) = f''(c)&\Longrightarrow 2a_2=f''(c) \Longrightarrow \boxed{a_2=\frac{f''(c)}{2}}\\
  P_n^{(3)}(c)=f^{(3)}(c) &\Longrightarrow 3\cdot2a_3=f^{(3)}(c) \Longrightarrow \boxed{a_3=\frac{f^{(3)}(c)}{3\cdot 2}}\\
  P_n^{(4)}(c)=f^{(4)}(c) &\Longrightarrow 4\cdot 3\cdot 2a_4=f^{(4)}(c)\Longrightarrow \boxed{a_4=\frac{f^{(4)}(c)}{4\cdot3\cdot 2}}
  \end{align*}
   Notice that we are purposefully leaving the denominators unsimplified in order to help us see a pattern.
 \end{solution}
 
    
  \item
  \begin{problem}[1in]
  Do you see a pattern? Can you write down a formula for $a_k$ (the coefficient of the $k$th degree term in the polynomial) in terms of $f^{(k)}(c)$?
  \end{problem}
  
  \begin{solution}
  Based on the pattern above, we see that the coefficient of the $k$th degree term for the Taylor polynomial is given by:
  $$\boxed{a_k = \frac{f^{(k)}(c)}{k!}}$$
  where $f^{(k)}(c)$ is the $k$th derivative of $f$ at $x=c$, and $k!=k\cdot (k-1)\cdot (k-2)\cdot~\cdots~\cdot 3\cdot 2 \cdot 1$. Note that we define $0!=1$, as it allows our formula to work for all $k$ values, including 0.
  
Notice that this formula is consistent with what we found in the previous parts. For example, $a_2=\frac{f^{(2)}(c)}{2!}=\frac{f''(c)}{2}$, $a_1=\frac{f^{(1)}(c)}{1!}=\frac{f'(c)}{1}=f'(c)$ and $a_0=\frac{f^{(0)}(c)}{0!}=\frac{f(c)}{1}=f(c)$. 
  
The boxed equation above is extremely important for you to know! It allows you to find the missing coefficients of the Taylor polynomial for a given function $f$, without needing to repeat the process we did in parts (a)-(c). From this point on, whenever you're asked to find the Taylor polynomial of a function, you can (and should!) use this formula without proof to find the coefficients, instead of setting derivatives equal to each other.
  \end{solution}

  \end{enumerate}

  
  

  
  
  

  
  
  \item 
  \note{Practice using the formula from {\#1}.  Students often get very confused by coefficients that are 0. It's
important to discuss why the odd coefficients are all 0 for $\cos x$. Make sure they also
  understand why the 15th degree polynomial for $\cos x$ ends with a degree 14 term. (A good question to ask here is ``what would the degree 14 polynomial for $\cos x$ be? What about the degree 13 polynomial?".) Also, emphasize that the number of terms in the polynomial is NOT the same thing as the degree of the polynomial.}

  \begin{problem}[2.25in]
  Find the degree $15$ Taylor polynomial approximation for $f(x) = \cos
  x$ centered at $0$.
  \end{problem}

  \begin{solution}
  The degree 15 Taylor polynomial approximation is $P_{15}(x) = a_0 +
  a_1 x + a_2 x^2 + \cdots + a_{15} x^{15}$, where $a_k =
  \frac{f^{(k)}(0)}{k!}$.  So, to find these coefficients, we need to
  compute derivatives of $f(x) = \cos x$:
  \[ \begin{array}{rcl c rcl}
  f(x) & = & \cos x & \Rightarrow & f(0) & = & 1 \\
  f'(x) & = & -\sin x & \Rightarrow & f'(0) & = & 0 \\
  f''(x) & = & -\cos x & \Rightarrow & f''(0) & = & -1 \\
  f'''(x) & = & \sin x & \Rightarrow & f'''(0) & = & 0
  \end{array}\]
  After this, the derivatives repeat in this pattern forever.  So,
  \begin{align*}
  a_0 &=\frac{f^{(0)}(0)}{0!}=\frac{1}{1}=1 \\
  a_1 &=\frac{f^{(1)}(0)}{1!}=\frac{0}{1}=0 \\
  a_2 &=\frac{f^{(2)}(0)}{2!}=\frac{-1}{2}= -\frac{1}{2!} \\
  a_3 &= \frac{f^{(3)}(0)}{3!}=\frac{0}{3!}=0 \\
  a_4 &= \frac{f^{(4)}(0)}{4!}=\frac{1}{4!} \\
  a_5 &= \frac{f^{(5)}(0)}{5!}=\frac{0}{5!}=0 \\
  &~~~~~~~~~\eqdots\\
   a_14 &= \frac{f^{(14)}(0)}{14!}=\frac{-1}{14!}=-\frac{1}{14!}\\
   a_15 &= \frac{f^{(15)}(0)}{15!}=\frac{0}{15!}=0\\
  \end{align*}
  Thus, 
  \begin{align*}
  P_{15}(x) &=a_0+a_1x+a_2x^2+a_3x^3+a_4x^4+\cdots + a_{14}x^{14} + a_{15}x^{15}\\
  &=1+(0)x+\left(-\frac{1}{2!}\right)x^2+(0)x^3+\left(\frac{1}{4!}\right)x^4+(0)x^5+\cdots + \left(-\frac{1}{14!}\right)x^{14}+(0)x^{15}\\
  &= \boxed{1 - \frac{x^2}{2!} + \frac{x^4}{4!} - \frac{x^6}{6!} + \frac{x^8}{8!} - \frac{x^{10}}{10!} +\frac{x^{12}}{12!} -\frac{x^{14}}{14!}}
  \end{align*}  
  Observe that we still refer to this polynomial as the degree 15 Taylor polynomial approximation, even though a 15th degree term does not appear in the polynomial (since its coefficient was 0). If we wanted the degree 14 Taylor polynomial approximation of $\cos x$ centered at 0, it would also be given by the same boxed polynomial above.
  
  Also note that although this is the degree 15 Taylor polynomial, it only has 8 terms (since the other 7 terms had coefficients of 0). This indicates that the \textit{degree} of a polynomial is different than the \textit{number of terms} in a polynomial.
  \end{solution}

   \pspace{\newpage}
  
  \item 
  Rewrite the following sums in $\cdots$
  notation.  (This means: Write down at least the first 3 terms $+ \cdots
  + $ the last term.)

  \begin{enumerate}
    \item

    \note{This shows them how to put a sum involving even numbers into
    sigma notation.  (It's worthwhile to point this out explicitly.)}
    
    \begin{problem}[\vfill]
    $\displaystyle \sum_{k=1}^{100} \frac{1}{2k}$
    \end{problem}

    \begin{solution}
    $\displaystyle \frac{1}{2} + \frac{1}{4} + \frac{1}{6} + \cdots +
    \frac{1}{198} + \frac{1}{200}$
    \end{solution}

    \item

    \note{This shows them how to put a sum whose terms alternate in sign
    into sigma notation.  I also like to ask how we would ``reverse''
    the signs (to write $\displaystyle \frac{1}{3} - \frac{1}{4} +
    \frac{1}{5} - \cdots + \frac{1}{15}$ instead).  If they can see that
    they simply need to multiply what they have by $-1$, then it's easy
    for them to get to changing $(-1)^k$ to $(-1)^{k+1}$.  I
    emphasize that, when we want to write terms that alternate in sign,
    we always use either $(-1)^k$ or $(-1)^{k+1}$.}
    
    \begin{problem}[\vfill]
    $\displaystyle \sum_{k=3}^{15} \frac{(-1)^k}{k}$
    \end{problem}

    \begin{solution}
    $\displaystyle -\frac{1}{3} + \frac{1}{4} - \frac{1}{5} + \cdots -
    \frac{1}{15}$
    \end{solution}

    \item

    \note{This shows them how to put odd numbers into sigma notation.}
    
    \begin{problem}[\vfill]
    $\displaystyle \sum_{k=0}^{25} \frac{x^{2k + 1}}{k + 1}$
    \end{problem}

    \begin{solution}
    $\displaystyle x + \frac{x^3}{2} + \frac{x^5}{3} + \frac{x^7}{4} +
    \cdots + \frac{x^{51}}{26}$
    \end{solution}

    \item
    
    \note{This is the same sum as the previous one.  Students sometimes
    try to give the index of summation a particular meaning, so I use
    this to point out that $k$ doesn't really mean anything on its own.}
    
    \begin{problem}[\vfill]
    $\displaystyle \sum_{k=1}^{26} \frac{x^{2k - 1}}{k}$
    \end{problem}

    \begin{solution}
    $\displaystyle x + \frac{x^3}{2} + \frac{x^5}{3} + \frac{x^7}{4} +
    \cdots + \frac{x^{51}}{26}$
    \end{solution}
    
  \end{enumerate}





  \item 
  Rewrite the following sums in sigma notation.
  \begin{enumerate}
    \item
    \begin{problem}[0.5in]
    $\displaystyle \frac{1}{3^3} + \frac{1}{4^3} + \frac{1}{5^3} +
    \frac{1}{6^3} + \cdots + \frac{1}{100^3}$
    \end{problem}

    \begin{solution}
    One possible solution is $\displaystyle \sum_{k=3}^{100}
    \frac{1}{k^3}$.  There are many possible solutions, depending on
    what you start $k$ at.  For instance, another valid answer is
    $\displaystyle \sum_{k=0}^{97} \frac{1}{(k + 3)^3}$
    \end{solution}

    \item

    \note{Help students with a strategy here.  I usually have them start
    by ignoring the fact that the terms alternate in sign.  We see that
    we're doing the square root of even numbers, so we start with
    $\sqrt{2k}$.  Then we figure out the starting and ending values of
    $k$, and finally we decide whether we should use $(-1)^k$ or
    $(-1)^{k+1}$ to make the terms alternate as we like (I emphasize
    that they should just check the first $k$-value to see which
    works).} 
    
    \begin{problem}[0.5in]
    $\displaystyle \sqrt{2} - \sqrt{4} + \sqrt{6} - \sqrt{8} + \sqrt{10}
    - \cdots - \sqrt{120}$
    \end{problem}

    \begin{solution}
    $\displaystyle \sum_{k=1}^{60} (-1)^{k + 1} \sqrt{2k}$ is a possible
    solution.
    \end{solution}

    \item
    \begin{problem}[0.5in]
    $\displaystyle 1 + x^3 + x^6 + x^9 + x^{12} + x^{15} + \cdots +
    x^{30}$
    \end{problem}

    \begin{solution}
    $\displaystyle \sum_{k=0}^{10} x^{3k}$ is a possible solution.
    \end{solution}

    \item
    \begin{problem}[0.5in]
    $x - x^3 + x^5 - x^7 + x^9 - \cdots + x^{101}$
    \end{problem}

    \begin{solution}
    $\displaystyle \sum_{k=0}^{50} (-1)^k x^{2k + 1}$ and $\displaystyle
    \sum_{k=1}^{51} (-1)^{k + 1} x^{2k - 1}$ are two possible solutions.
    \end{solution}
    
    \item
    \begin{problem}[0.5in]
    $\displaystyle x - \frac{x^3}{2} + \frac{x^5}{3} - \frac{x^7}{4} +
    \cdots - \frac{x^{31}}{16}$
    \end{problem}
    
    \begin{solution}
    $\displaystyle \sum_{k=0}^{15} (-1)^k \frac{x^{2k + 1}}{k + 1}$ and
    $\displaystyle \sum_{k=1}^{16} (-1)^{k + 1} \frac{x^{2k - 1}}{k}$
    are two possible solutions.
    \end{solution}
  \end{enumerate}

\pspace{\newpage}

  \item 
  \begin{enumerate}
    \item 
    \begin{problem}[1in]
    Rewrite the degree 15 Taylor polynomial for $\cos(x)$ centered at 0 in
    sigma notation.
    \end{problem}
    
    \begin{solution}
    $\displaystyle \sum_{k=0}^7 (-1)^k \frac{x^{2k}}{(2k)!}$ is one
    possibility. 
    \end{solution}

    \item
    \begin{problem}[1.5in]
    Find the degree 25 Taylor polynomial approximation for $f(x) = \cos
    x$ centered at 0.  Write it in both sigma and $\cdots$ notation.
    \end{problem}

    \begin{solution}
    We already found the pattern for the coefficients in
    {\#2}; we just need to add all the terms
    up to degree 25 instead of 15.  Since we want degree 25, the last
    term will be $\frac{x^{24}}{24!}$, which corresponds to $k = 12$
    (the coefficient of the $x^{25}$ term is 0).  So, the degree 25
    Taylor polynomial approximation is \fbox{$\displaystyle
    \sum_{k=0}^{12} (-1)^k \frac{x^{2k}}{(2k)!} = 1 - \frac{x^2}{2!} +
    \frac{x^4}{4!} - \cdots + \frac{x^{24}}{24!}$}.
    \end{solution}

    \item

    \note{If you want to use Mathematica to look at this, be aware that
    Mathematica considers all decimals to have very low precision.
    You'll get better results if you enter 0.1 as the fraction 1/10 than
    if you just put in 0.1.}
    
    \begin{problem}[1.5in]
    Use the degree 25 Taylor polynomial approximation to get an
    approximation for $\cos 0.1$. You may leave your answer as a sum of fractions.
    \end{problem}

    \begin{solution}
    We simply plug $x = 0.1$ into our degree 25 Taylor polynomial
    approximation to get \fbox{$\displaystyle \sum_{k=0}^{12} (-1)^k
    \frac{0.1^{2k}}{(2k)!} = 1 - \frac{0.1^2}{2!} + \frac{0.1^4}{4!} -
    \cdots + \frac{0.1^{24}}{24!}$}.

    (If you add this up, you get
    $\frac{617348744064667103320115839788837194165430184801}{620448401733239439
    360000000000000000000000000000}$, which differs from $\cos 0.1$ by
    about $2.5 \times 10^{-53}$ --- not bad!)
    \end{solution}
  \end{enumerate}

  \begin{framed}
The \textbf{Taylor series} of $f(x)$ centered at $x=c$ is given by:
$$\sum_{k=0}^{\infty} a_k(x-c)^k =a_0+a_1(x-c)+a_2(x-c)^2+ \cdots$$
where $$a_k=\frac{f^{(k)}(c)}{k!}$$
\end{framed}
\begin{framed}
  \centerline{\textbf{The Main Questions of the Series Unit}}
  We hope that, by using a ``polynomial of infinite degree,'' we end up with
  something that is not just an approximation for our function but is
  actually equal to the function. In order to know if/when this is actually true, we need to answer two main questions:
  \begin{enumerate}
  \item \textbf{For which functions $f$ and for which inputs $x$ does the Taylor series of $f(x)$ centered at $c$ actually add up to anything meaningful?} For example, does the Taylor series for $\sin x$, given by $x - \frac{x^3}{3!} + \frac{x^5}{5!} - \cdots$, even have a defined value when $x=1$? When $x=-1$? When $x=100$?
  \item \textbf{For which functions $f$ and for which inputs $x$ does the Taylor series of $f(x)$ centered at $c$ actually add up \textit{specifically to} $f(x)$?} For example, perhaps the Taylor series $x - \frac{x^3}{3!} + \frac{x^5}{5!} - \cdots$ has a defined value when $x=100$, but is \textit{not} actually equal to $\sin(100)$. (In this case, our Taylor series for $\sin x$ wouldn't be very useful!)
  \end{enumerate}
  \end{framed}

\end{enumerate}


\end{document}